% Generated from the audited Markdown source; author: Carptopus.
\documentclass[11pt]{article}
\usepackage[a4paper,margin=27mm]{geometry}
\usepackage[T1]{fontenc}
\usepackage[utf8]{inputenc}
\usepackage{lmodern}
\usepackage{microtype}
\usepackage{amsmath,amssymb,amsthm,mathtools}
\usepackage{needspace}
\usepackage{xcolor}
\usepackage[colorlinks=true,linkcolor=blue!55!black,citecolor=blue!55!black,urlcolor=blue!65!black]{hyperref}
\usepackage{fancyhdr}

\newtheorem{theorem}{Theorem}[section]
\newtheorem{lemma}[theorem]{Lemma}
\newtheorem{proposition}[theorem]{Proposition}
\newtheorem{corollary}[theorem]{Corollary}
\theoremstyle{remark}
\newtheorem{remark}[theorem]{Remark}
\numberwithin{equation}{section}
\allowdisplaybreaks[2]

\pagestyle{fancy}
\fancyhf{}
\fancyhead[L]{Carptopus}
\fancyhead[R]{Fano half-rank profiles}
\fancyfoot[C]{\thepage}
\setlength{\headheight}{14pt}
\setlength{\parindent}{0pt}
\setlength{\parskip}{0.45em}
\setlength{\emergencystretch}{4em}
\urlstyle{same}

\title{Realizability and minimum dimension\\
of Fano half-rank profiles over \texorpdfstring{$\mathbb F_2$}{F2}}
\author{Carptopus\\\href{mailto:carptopus@163.com}{carptopus@163.com}}
\date{30 August 2026}

\hypersetup{
  pdftitle={Realizability and minimum dimension of Fano half-rank profiles over F2},
  pdfauthor={Carptopus},
  pdfsubject={A classification of realizable Fano half-rank profiles and their minimum ambient dimensions over F2},
  pdfkeywords={alternating bilinear form, alternating matrix space, Fano plane, rank profile, finite field, Hilbert basis, minimum dimension, Reed-Muller code}
}

\begin{document}
\maketitle

\begin{center}
\fcolorbox{black!35}{black!4}{\parbox{0.88\textwidth}{\centering\small
Public Beta manuscript, version 0.1-beta. The mathematical proof and the
complete integrated manuscript have passed project-internal zero-trust audits.
External mathematical review and formal peer review are pending.}}
\end{center}

\begin{abstract}


Let $V$ be a finite-dimensional vector space over $\mathbb F_2$ and let

\nopagebreak[4]
\[
B:\mathbb F_2^3\longrightarrow \mathrm{Alt}(V)
\]

be a linear family of alternating bilinear forms. Its half-rank profile is the
integer-valued function

\nopagebreak[4]
\[
h(a)=\frac{1}{2}\mathrm{rank}(B_a),
\qquad a\in\mathbb F_2^3\setminus\{0\}.
\]

We classify these labelled seven-point profiles in two senses. First, we
determine exactly which nonnegative integer functions on the Fano plane occur in
some finite dimension. The triangle inequalities are sufficient except for
seven primitive line jumps and seven infinite parity cosets supported on
singleton faces. Second, for every realizable profile $h$ we determine the
minimum possible dimension $\mu(h)$ of $V$. If $t=\max h$, then
$\mu(h)\in\{2t,2t+1\}$, and we give a complete description of the profiles that
require the extra dimension: two height-one orbits, three height-two orbits,
and one explicit infinite family $E_t$.

The nonrealizability arguments use Pfaffian parity, a simultaneous decomposition
of commuting self-adjoint idempotents, and a rank-two Pfaffian update identity.
The sufficiency proofs are computer assisted only at finite, exact interfaces:
an integral Hilbert-basis certificate and a conductor reduction to a verified
collection of profiles of height at most six.

\end{abstract}

\section{Introduction}


Rank distributions of linear spaces of alternating matrices arise in finite
geometry, coding theory, and the study of bilinear and quadratic forms. The
two-dimensional parameter case asks for the ranks of a triple
$(A,B,A+B)$. It has been studied in a precise enumerative form, for example by
Pott, Schmidt, and Zhou \cite{pott2016}. Passing from a two-dimensional parameter space to a
three-dimensional one replaces the three nonzero parameters by the seven points
of the Fano plane. The resulting seven ranks are not independent: each Fano line
is a two-dimensional subfamily and therefore satisfies rank subadditivity.

The purpose of this paper is to answer two natural questions for this first
seven-point case over $\mathbb F_2$.

\begin{enumerate}
\item Which labelled seven-tuples occur as the half-ranks of a linear net of
   alternating forms in some finite dimension?
\item For a realizable profile, what is the least dimension in which it occurs?

\end{enumerate}

The answer to the first question is almost, but not quite, the integral metric
cone of the Fano plane. There are two kinds of holes. One is finite: a primitive
profile that jumps from one off a line to two on the line. The other consists of
seven infinite parity cosets on four-generated faces of the cone. Both
obstructions have structural explanations rather than being artifacts of a
bounded search.

The answer to the second question is unexpectedly rigid. The evident lower
bound is twice the largest half-rank. Every realizable profile meets this bound
or misses it by exactly one. The profiles that miss it form a short finite list
in heights one and two and a single infinite family obtained by repeatedly
adding one Fano cut.

Our classification is a classification of labelled rank profiles, not of
alternating-matrix spaces up to isomorphism. Eisenbud and Koh classified
three-dimensional spaces of $4\times4$ and $5\times5$ alternating matrices over
an algebraically closed field \cite{eisenbud1994}. Gow obtained dimension bounds for spaces with
restricted rank sets over finite fields \cite{gow2017}, while Qiao studied the enumeration
of alternating-matrix spaces with explicit coordinates \cite{qiao2021}. These results are
closely related in language but have a different output. Here the field is fixed
to $\mathbb F_2$, the matrix size is arbitrary, and the output is a complete
seven-point realizability and minimum-dimension criterion.

The proof has two finite exact components. We certify the Hilbert basis of the
Fano triangle cone and later certify a finite conductor for sharp-dimensional
realizability. Both certificates use integer or binary arithmetic; their role
and logical boundary are stated explicitly in Sections 3, 7, and 8.

\section{Definitions and basic properties}


Throughout, $V$ is a finite-dimensional vector space over $\mathbb F_2$.
Write $\mathrm{Alt}(V)$ for the vector space of alternating bilinear forms on
$V$. Equivalently, after choosing a basis, its elements are symmetric matrices
with zero diagonal. Their ranks are even. We use
$\mathbb N=\{0,1,2,\ldots\}$.

Let

\nopagebreak[4]
\[
P=\mathbb F_2^3\setminus\{0\}.
\]

The seven elements of $P$ are the labelled points of the Fano plane. Its lines
are the triples

\nopagebreak[4]
\[
\{a,b,a+b\},\qquad a,b\in P,\ a\ne b.
\]

An alternating net is a linear map

\nopagebreak[4]
\[
B:\mathbb F_2^3\longrightarrow\mathrm{Alt}(V).
\]

Its half-rank profile is the function $h_B:P\to\mathbb N$ given by

\nopagebreak[4]
\[
h_B(a)=\frac{1}{2}\mathrm{rank}(B_a).
\tag{2.1}
\]

We call $h$ realizable if $h=h_B$ for some $V$ and some alternating net $B$.
The realizable profiles form an additive semigroup: orthogonal direct sum of
nets adds profiles coordinatewise.

For a Fano line $L$ define the cut profile $c_L$ and the line jump $\ell_L$ by

\nopagebreak[4]
\[
c_L(x)=
\begin{cases}
0,&x\in L,\\
1,&x\notin L,
\end{cases}
\qquad
\ell_L(x)=
\begin{cases}
2,&x\in L,\\
1,&x\notin L.
\end{cases}
\tag{2.2}
\]

For $p\in P$, define the singleton jump

\nopagebreak[4]
\[
s_p(x)=
\begin{cases}
2,&x=p,\\
1,&x\ne p.
\end{cases}
\tag{2.3}
\]

Rank subadditivity gives, on every Fano line,

\nopagebreak[4]
\[
h(a+b)\le h(a)+h(b).
\tag{2.4}
\]

Let $C_{\mathbb Z}$ be the semigroup of all nonnegative integer functions on
$P$ satisfying the three versions of (2.4) on each line. We call it the
integral Fano triangle cone.

The effective dimension of a realizable profile is

\nopagebreak[4]
\[
\mu(h)=\min\{\dim V:h=h_B\text{ for a net on }V\}.
\tag{2.5}
\]

Adding a common radical does not change the profile. Consequently, once $h$ is
realizable in dimension $\mu(h)$, it is realizable in every dimension at least
$\mu(h)$.

\section{The integral Fano triangle cone}


The cone $C_{\mathbb Z}$ provides the necessary inequalities and the finite
additive framework used in the realizability proof.

\Needspace{0.28\textheight}
\begin{lemma}[extreme rays and Hilbert basis]

The rational cone underlying $C_{\mathbb Z}$ has fourteen primitive extreme
rays: the seven cut profiles $c_L$ and the seven line jumps $\ell_L$.
Its integral Hilbert basis has 78 elements. Of these, 64 are realizable in
dimension at most four. The remaining fourteen are precisely the seven
$s_p$ and the seven $\ell_L$.

\subsubsection*{Exact certificate}


There are 28 defining inequalities: seven nonnegativity inequalities and three
triangle inequalities on each of seven lines. The program
\nolinkurl{verify_fano_triangle_extreme_rays.cpp} examines every six-subset of the 28
inequalities, computes its generalized cross product over the integers, and
retains the primitive feasible rays. It finds exactly the fourteen profiles
above.

The program \nolinkurl{verify_fano_triangle_hilbert_basis.py} reconstructs the facet
incidences, forms an exact pulling triangulation with twelve simplicial cones,
and enumerates their half-open fundamental parallelepipeds. Every determinant,
adjugate, divisibility test, and decomposability test is exact. The union gives
the stated 78-element Hilbert basis. The alternating-net witnesses for the
realizable atoms are reconstructed and checked over $\mathbb F_2$ by
\nolinkurl{verify_realized_strengthened_fano_atoms.py}.

The lemma is therefore a finite computer-assisted lemma. It does not infer a
general statement from a numerical bound.

\end{lemma}

\Needspace{0.28\textheight}
\begin{lemma}[multiples of missing atoms]

For every point $p$ and every line $L$,

\nopagebreak[4]
\[
2s_p=\sum_{K:p\notin K}c_K.
\tag{3.1}
\]

Moreover, $m\ell_L$ is realizable for every $m\ge2$.

\end{lemma}

\begin{proof}

There are four Fano lines avoiding $p$, and direct evaluation gives (3.1).
Each $c_K$ is realized on a two-dimensional symplectic space, so $2s_p$ is
realizable by orthogonal direct sum.

For the line jump, explicit matrix spaces give

\nopagebreak[4]
\[
2\ell_L=(4,4,4,2,2,2,2),
\qquad
3\ell_L=(6,6,6,3,3,3,3),
\tag{3.2}
\]

after a suitable labelling with $L$ in the first three coordinates. The
matrices and all seven ranks are checked in
\nolinkurl{verify_fano_43_counterexample.py}. Every integer $m\ge2$ is a nonnegative
combination of 2 and 3, so orthogonal direct sums finish the proof.
\end{proof}


\section{The realizability semigroup}


We now identify the two genuine boundaries left by Lemma 3.1.

\Needspace{0.28\textheight}
\begin{lemma}[rank-additive reduction]

Let $A,C:V\to V^*$ be linear maps. If

\nopagebreak[4]
\[
\mathrm{rank}(A+C)=\mathrm{rank}(A)+\mathrm{rank}(C),
\tag{4.1}
\]

then

\nopagebreak[4]
\[
\mathrm{im}(A)\cap\mathrm{im}(C)=0,
\qquad
\ker(A+C)=\ker(A)\cap\ker(C).
\tag{4.2}
\]

\end{lemma}

\begin{proof}

The image of $A+C$ is contained in $\mathrm{im}(A)+\mathrm{im}(C)$.
Equality of dimensions in (4.1) forces the latter sum to be direct and forces
$\mathrm{im}(A+C)$ to equal it. If $(A+C)x=0$, then $Ax=Cx$ lies in the
intersection of the two images and hence both terms vanish. The reverse kernel
inclusion is immediate.
\end{proof}


\Needspace{0.28\textheight}
\begin{proposition}[primitive line jumps are not realizable]

For every Fano line $L$, the profile $\ell_L$ is not realizable.

\end{proposition}

\begin{proof}

Write $L=\{u,v,u+v\}$ and choose $w\notin L$. Suppose $\ell_L$ were realized.
Then $P=B_u$ has rank four, while

\nopagebreak[4]
\[
B_w,\ B_{w+u},\ B_{w+v},\ B_{w+u+v}
\]

all have rank two. Both decompositions

\nopagebreak[4]
\[
P=B_w+B_{w+u}=B_{w+v}+B_{w+u+v}
\tag{4.3}
\]

are rank additive. Lemma 4.1 shows that $\mathrm{rad}(P)$ lies in the radical
of each of the four rank-two forms. The other two forms on $L$ are sums of
these forms, so the entire net descends to $V/\mathrm{rad}(P)$, which has
dimension four.

On a four-dimensional space an alternating form is nondegenerate exactly when
its maximal Pfaffian is nonzero. The Pfaffian of the net is a zero-constant
Boolean polynomial of degree at most two in the three parameters. Every such
function has even Hamming weight on $\mathbb F_2^3$: each monomial of degree
less than three has even total weight, and the value at zero is zero. The
profile $\ell_L$, however, requires the Pfaffian to be nonzero exactly on the
three points of $L$. This odd weight is impossible.
\end{proof}


\Needspace{0.28\textheight}
\begin{proposition}[the singleton face]

Fix $p\in P$. Suppose a realizable profile $h$ satisfies equality in all three
triangle inequalities whose opposite point is $p$; equivalently,

\nopagebreak[4]
\[
h(p)=h(a)+h(a+p)
\tag{4.4}
\]

for the three unordered pairs $\{a,a+p\}$ in $P\setminus\{p\}$. Then

\nopagebreak[4]
\[
h\in\mathbb N\{c_L:p\notin L\}.
\tag{4.5}
\]

Conversely, every profile in the semigroup on the right is realizable and
satisfies (4.4).

\end{proposition}

\begin{proof}

Let $P_0=B_p$. Lemma 4.1 applied to each equality in (4.4) shows that all forms
descend to $V/\mathrm{rad}(P_0)$. We may therefore suppose that $P_0$ is
nondegenerate.

Choose two other parameter directions and write the corresponding forms as

\nopagebreak[4]
\[
A=P_0T,
\qquad
C=P_0S.
\]

We use the following elementary observation. If an endomorphism $T$ satisfies

\nopagebreak[4]
\[
\mathrm{rank}(T)+\mathrm{rank}(I+T)=\dim V,
\tag{4.6}
\]

then $T$ is idempotent. Indeed,
$\mathrm{im}(T)+\mathrm{im}(I+T)=V$, so (4.6) makes the sum direct. The vector
$(I+T)Tx=T(I+T)x$ lies in both summands and must vanish.

The three equalities (4.4) imply that $T$, $S$, and $T+S$ satisfy (4.6), hence
are idempotent. Alternation of $A$ and $C$ says that $T$ and $S$ are
self-adjoint with respect to $P_0$. In characteristic two,

\nopagebreak[4]
\[
(T+S)^2=T+S
\]

together with $T^2=T$ and $S^2=S$ gives $TS=ST$.

Thus $V$ is the direct sum of the four simultaneous eigenspaces $V_{ij}$ for
$(T,S)$. Self-adjointness makes distinct eigenspaces orthogonal for $P_0$.
Since $P_0$ is nondegenerate, its restriction to each $V_{ij}$ is
nondegenerate; consequently every $V_{ij}$ has even dimension.

On $V_{ij}$ each parameter form is a scalar linear function on
$\mathbb F_2^3$ times the fixed symplectic form $P_0|_{V_{ij}}$. That scalar is
one at $p$, so its zero set is a line avoiding $p$. If
$\dim V_{ij}=2m_{ij}$, this block contributes $m_{ij}$ times the corresponding
cut profile. Summing the four blocks proves (4.5). The converse follows from
the standard two-dimensional realizations of the cuts.
\end{proof}


\Needspace{0.28\textheight}
\begin{corollary}[singleton parity cosets are holes]

For every $p\in P$, no profile in

\nopagebreak[4]
\[
s_p+\mathbb N\{c_L:p\notin L\}
\tag{4.7}
\]

is realizable.

\end{corollary}

\begin{proof}

All profiles in (4.7) satisfy the equalities (4.4). The four cut vectors in
(4.5) are linearly independent on that face. Identity (3.1) says that $s_p$ is
the nontrivial order-two lattice coset represented by half their sum. Hence the
coset (4.7) is disjoint from the semigroup (4.5), while Proposition 4.3 places
every realizable profile on the face in (4.5).
\end{proof}


The remaining sufficiency step is finite once the Hilbert basis is known.

\Needspace{0.28\textheight}
\begin{lemma}[finite coupling certificate]

Let $H$ be the 78-element Hilbert basis from Lemma 3.1, and let

\nopagebreak[4]
\[
M=\{s_p:p\in P\}\cup\{\ell_L:L\text{ a Fano line}\}
\]

be its fourteen nonrealizable atoms. There is a finite exact witness library
with the following properties.

\begin{enumerate}
\item Every sum of two atoms in $M$ is realizable.
\item Every sum of three atoms in $M$ is realizable, except $3s_p$.
\item The sum of any $\ell_L$ and any atom of $H\setminus M$ is realizable.
\item The sum $s_p+g$, for $g\in H\setminus M$, is nonrealizable exactly when
   $g=c_L$ for a line $L$ avoiding $p$.

\end{enumerate}

\subsubsection*{Exact certificate}


The program \nolinkurl{verify_fano_profile_two_saturation.py} reconstructs 218 profiles
with separately checked alternating-net witnesses, closes them additively in
the finite target box, and verifies all 105 unordered pairs, 560 unordered
triples, and 896 mixed targets. It reports precisely the four rules above. All
profile additions and rank checks are exact.

\end{lemma}

\Needspace{0.28\textheight}
\begin{theorem}[realizability classification]

A profile $h:P\to\mathbb N$ is realizable in some finite dimension if and only
if $h\in C_{\mathbb Z}$ and neither of the following holds:

\begin{enumerate}
\item $h=\ell_L$ for a Fano line $L$;
\item for some $p\in P$,

\nopagebreak[4]
\[
   h\in s_p+\mathbb N\{c_L:p\notin L\}.
   \tag{4.8}
\]

\end{enumerate}

\end{theorem}

\begin{proof}

Necessity of the triangle inequalities follows from rank subadditivity.
Proposition 4.2 and Corollary 4.4 prove necessity of the two exclusions.

Conversely, decompose $h$ into elements of the Hilbert basis $H$. The elements
of $H\setminus M$ are realizable. It remains to absorb the multiset of missing
atoms.

If its size is even, pair the missing atoms and use Lemma 4.5(1). If its size is
odd and at least three, a mixed triple can first be realized by Lemma 4.5(2),
after which the rest can be paired. The only unmixed cases require separate
attention. Repeated line jumps are realizable from multiplicity two onward by
Lemma 3.2. Repeated singleton jumps are realizable in even multiplicity by
(3.1), while an odd multiplicity reduces to one $s_p$ plus cuts avoiding $p$;
this is exactly the excluded coset (4.8).

Finally, if one singleton atom remains together with good Hilbert atoms,
Lemma 4.5(4) says that failure occurs exactly when all additions preserve the
same singleton face, again giving (4.8). A remaining line jump is absorbed by
Lemma 4.5(3), unless it is the primitive line jump itself. Thus every profile
not in the two excluded families is realizable.
\end{proof}


\subsection*{Regression check}


As a destructive check, \nolinkurl{probe_fano_profile_semigroup_holes.py} enumerates all
61,075 points of $C_{\mathbb Z}$ with $\max h\le6$. The witness semigroup
contains 60,578 of them. The remaining 497 profiles agree point for point with
the two families in Theorem 4.6. This bounded enumeration is not used to infer
the theorem.

\section{Minimum effective dimension}


For a nonzero realizable profile $h$, set

\nopagebreak[4]
\[
t=\max_{a\in P}h(a).
\tag{5.1}
\]

Every realization has dimension at least $2t$. We now classify exactly when
one additional dimension is necessary.

Choose linearly independent points $a,b,c\in P$, and put
$s=a+b+c$. For $t\ge3$, define $E_{t;a,b,c}$ by

\nopagebreak[4]
\[
E_{t;a,b,c}(x)=
\begin{cases}
t,&x\in\{a,b,c\},\\
t-1,&x=s,\\
1,&x\in\{a+b,a+c,b+c\}.
\end{cases}
\tag{5.2}
\]

The group $\mathrm{GL}(3,2)$ has one orbit of these profiles, of labelled size
28 for each $t$.

\Needspace{0.28\textheight}
\begin{theorem}[minimum-dimension classification]

Let $h$ be a nonzero realizable profile and let $t=\max h$. Then

\nopagebreak[4]
\[
\mu(h)=2t+\varepsilon(h),
\qquad \varepsilon(h)\in\{0,1\}.
\tag{5.3}
\]

The value $\varepsilon(h)=1$ occurs exactly in the following cases.

\begin{enumerate}
\item \textbf{Height one.} The profile is constant one, or it has one zero coordinate
   and six coordinates equal to one.
\item \textbf{Height two.} Every coordinate is positive, and the set
   $\{x:h(x)=2\}$ is a noncollinear three-set, the complement of a two-set, or
   all seven points.
\item \textbf{The infinite family.} The profile is $E_{t;a,b,c}$ for some $t\ge3$ and
   some linearly independent $a,b,c$.

\end{enumerate}

The zero profile satisfies $\mu(0)=0$.

We prove necessity of the three exceptional families first.

\end{theorem}

\Needspace{0.28\textheight}
\begin{proposition}[height one]

The two height-one profiles in Theorem 5.1 have minimum dimension three. Every
other realizable height-one profile is a cut and has minimum dimension two.

\end{proposition}

\begin{proof}

On a two-dimensional space, $\mathrm{Alt}(V)$ is one-dimensional. The set of
parameters giving the unique nonzero form is the complement of the kernel of a
linear functional, hence is the complement of a Fano line. These are exactly
the cuts.

On a three-dimensional space every nonzero alternating form has rank two.
A rank-two linear map from $\mathbb F_2^3$ to $\mathrm{Alt}(\mathbb F_2^3)$
with a prescribed one-dimensional kernel gives the single-zero profile, while
an isomorphism gives the constant-one profile.
\end{proof}


\Needspace{0.28\textheight}
\begin{proposition}[height two]

The three height-two orbits in Theorem 5.1 have minimum dimension five.

\end{proposition}

\begin{proof}

In dimension four, the set of full-rank directions is the support on $P$ of the
maximal Pfaffian, a zero-constant Boolean polynomial of degree at most two.
As in Proposition 4.2, its weight is even. The three stated orbits have,
respectively, 3, 5, and 7 full-rank directions, so none occurs in dimension
four.

Exact five-dimensional alternating-net witnesses exist for every member of
the three orbits. They are included in the verified low-layer construction
library; \nolinkurl{probe_n5_profile_saturation.cpp} reconstructs the complete target set
and checks the ranks over $\mathbb F_2$. Hence the minimum is five.
\end{proof}


\Needspace{0.28\textheight}
\begin{lemma}[rank-two Pfaffian update]

Let $A$ and $X$ be alternating forms on a $2t$-dimensional space. For fixed
$A$, there is a linear functional $\lambda_A$ on $\mathrm{Alt}(V)$ such that

\nopagebreak[4]
\[
\mathrm{Pf}(A+X)=\mathrm{Pf}(A)+\lambda_A(X)
\tag{5.4}
\]

whenever $\mathrm{rank}(X)\le2$.

\end{lemma}

\begin{proof}

The principal-minor expansion of the Pfaffian is

\nopagebreak[4]
\[
\mathrm{Pf}(A+X)
=\sum_{I\subseteq[2t],\ |I|\text{ even}}
 \mathrm{Pf}(X_I)\mathrm{Pf}(A_{I^c}).
\tag{5.5}
\]

If $\mathrm{rank}(X)\le2$, then every principal submatrix $X_I$ of order at
least four is singular, so its Pfaffian vanishes. Only the empty set and the
two-element subsets remain. The empty set gives $\mathrm{Pf}(A)$, and the
two-element terms are linear in the entries of $X$. Their sum is a single
linear functional $\lambda_A$ that works simultaneously for every rank-two
$X$.
\end{proof}


\Needspace{0.28\textheight}
\begin{proposition}[the family $E_t$]

For every $t\ge3$,

\nopagebreak[4]
\[
\mu(E_{t;a,b,c})=2t+1.
\tag{5.6}
\]

\end{proposition}

\begin{proof}

Suppose first that $E_t$ were realized in dimension $2t$. Use a full-rank
direction as an origin and write the forms in three full-rank directions as

\nopagebreak[4]
\[
A,\qquad A+R,\qquad A+Q.
\]

The three pair-sum directions in (5.2) have rank two, so
$R,Q,R+Q$ all have rank two. Over $\mathbb F_2$, a nondegenerate alternating
form has Pfaffian one. Lemma 5.4 and nondegeneracy of $A+R$ and $A+Q$ imply

\nopagebreak[4]
\[
\lambda_A(R)=\lambda_A(Q)=0.
\]

Linearity gives $\lambda_A(R+Q)=0$, and hence $A+R+Q$ is also nondegenerate.
But the direction $s$ in (5.2) is required to have rank $2t-2$. This is a
contradiction, so $\mu(E_t)\ge2t+1$.

An exact seven-dimensional witness realizes $E_3$. Let

\nopagebreak[4]
\[
L=\{a+b,a+c,b+c\}.
\]

Then coordinatewise

\nopagebreak[4]
\[
E_{t+1}=E_t+c_L.
\tag{5.7}
\]

Taking an orthogonal direct sum with the two-dimensional realization of $c_L$
therefore realizes $E_t$ in dimension $2t+1$ for every $t\ge3$.
\end{proof}


\section{Sharp realizations and the conductor theorem}


Call a realizable profile $h$ \textbf{sharp} if it is realizable in dimension
$2\max h$; in particular, the zero profile is sharp in dimension zero.
Propositions 5.2, 5.3, and 5.5 show that the listed exceptional
profiles are not sharp. To prove that there are no others, we reduce every
height to a finite exact calculation.

By the action of $\mathrm{GL}(3,2)$, choose an anchor point labelled $1$ at
which $h$ attains its maximum. Consider the chamber

\nopagebreak[4]
\[
C_1=\{h\in C_{\mathbb Z}:h(1)\ge h(x)\text{ for all }x\in P\}.
\tag{6.1}
\]

\Needspace{0.28\textheight}
\begin{lemma}[sharp chamber generators]

The cone $C_1$ has 30 primitive rays. Replace them by the following scaled
profiles:

\begin{itemize}
\item leave each cut ray unchanged;
\item multiply each of the six anchored single-zero height-one rays by two;
\item multiply the constant height-one ray by three;
\item multiply each primitive line-jump ray by two;
\item leave every other primitive ray unchanged.

\end{itemize}

Every resulting profile is sharp, and all have the same anchor as a maximum
coordinate. Consequently, every nonnegative integer sum of them is sharp.

\end{lemma}

\begin{proof}

The cut constructions are two-dimensional. For the two exceptional scaled
height-one rays, explicit sharp witnesses are required rather than direct sums
of the three-dimensional realizations in Proposition 5.2. In dimension four,
the alternating edge-mask triple $(12,0,22)$ has profile

\nopagebreak[4]
\[
(2,0,2,2,2,2,2),
\]

and Fano relabelling gives all doubled single-zero rays. In dimension six, let
$K=\mathbb F_8$ and put $V=K\oplus K$. The trace-pairing forms

\nopagebreak[4]
\[
B_a((x,y),(x',y'))=\mathrm{Tr}_{K/\mathbb F_2}
\bigl(a(xy'+x'y)\bigr)
\]

are nondegenerate for every $a\ne0$. Restricting $a$ to a three-dimensional
$\mathbb F_2$-basis of $K$ gives the constant profile $(3,3,3,3,3,3,3)$.
Both constructions are reconstructed by
\nolinkurl{verify_scaled_chamber_ray_witnesses.py}.

Lemma 3.2 gives sharp realizations of doubled line jumps. The remaining rays
have exact sharp witnesses in the low-layer library.

For orthogonal direct sums, profiles add and dimensions add. Since the chosen
anchor is maximal in every summand, its coordinate equals the sum of the
individual maxima. Thus the resulting dimension is twice the maximum of the
sum profile.
\end{proof}


The factor three for the constant ray is essential: the constant-two profile is
one of the height-two exceptions.

\Needspace{0.28\textheight}
\begin{lemma}[finite conductor certificate]

Every lattice point of $C_1$ is exactly one of the following:

\begin{enumerate}
\item a sharp profile;
\item a nonrealizable profile excluded by Theorem 4.6;
\item one of the height-one or height-two exceptions of Theorem 5.1;
\item a member of the family $E_t$.

\end{enumerate}

\subsubsection*{Exact certificate and reduction}


The completeness of the 30 primitive chamber rays is independently certified
before the conductor calculation. The program
\nolinkurl{verify_fano_anchor_cone_extreme_rays.cpp} constructs all 34 inequalities
(nonnegativity, Fano triangles, and six anchor inequalities), examines all
1,344,904 six-subsets, and proves that their primitive feasible null rays are
exactly the 30 rays used by the conductor.

The 30 scaled rays in Lemma 6.1 then have a pulling triangulation into 81
full-dimensional simplicial cones. Their determinant distribution is

\nopagebreak[4]
\[
\{8:24,\ 16:36,\ 24:18,\ 48:3\}.
\tag{6.2}
\]

The union of their half-open fundamental parallelepipeds contains 226 integer
residue profiles, all of height at most six. In any fixed simplicial cone, every
lattice point has a unique decomposition

\nopagebreak[4]
\[
h=r+\sum_i m_i g_i,
\qquad m_i\in\mathbb N,
\tag{6.3}
\]

where $r$ is the residue and the $g_i$ are the seven scaled rays of that cone.

The exact transition table has five kinds of residues. A sharp residue remains
sharp after additions. A height-two exception is exceptional only with no
addition; the first addition reaches a verified sharp base. A height-one
exception either becomes sharp, passes once through a height-two exception, or
follows one repeated cut and generates $E_3,E_4,\ldots$. A nonrealizable
singleton-face residue either stays in its forbidden coset or leaves the face
and reaches a sharp base. The zero residue is generated directly by the rays.

All resulting sharp bases have height at most six. There are 812 distinct such
bases. \nolinkurl{verify_fano_graded_conductor.py} certifies the 81 cones, 226 residues,
and every transition using exact integer arithmetic. The normalized list of 812
bases is frozen in \nolinkurl{verification/results/required_sharp_bases.tsv};
\nolinkurl{verify_fano_graded_required_bases.py} regenerates the conductor output and
checks exact set equality with that file.

The witness side is fail closed. Complete five- and six-dimensional low layers
are reconstructed by \nolinkurl{probe_n5_profile_saturation.cpp} and
\nolinkurl{probe_n6_profile_saturation.cpp}; the exceptional $E_3$ witness is checked by
\nolinkurl{probe_n7_noncollinear_full_exception.cpp}. The dimension-eight boundary is
constructed by exact binary-rank searches plus a fixed targeted witness. The
remaining height-five and height-six representatives are checked by the
finite-field trace and explicit matrix scripts in the verification package.
Every search uses a fixed seed and exits nonzero unless its entire stated target
set is covered.

Finally, \nolinkurl{verify_fano_graded_excess_bound.cpp} reconstructs the complete
low-height chamber from these witness-certified generators and checks every one
of the frozen 812 bases at sharp dimension. It reports

\begin{quote}\small\raggedright
\texttt{REFINED\_\allowbreak{}CLASSIFICATION errors=0}\par
\texttt{REQUIRED\_\allowbreak{}BASE\_\allowbreak{}COVER bases=812 errors=0}\par
\end{quote}


Randomized routines were used only to discover some matrices. They have fixed
seeds, and every retained matrix is checked by exact Gaussian elimination. A
failure to find a matrix is never used as a proof of nonrealizability.

\end{lemma}

\begin{proof}[Proof of Theorem 5.1]

The lower bound $\mu(h)\ge2t$ follows from the definition of $t$. The zero
profile is realized on the zero space. Propositions 5.2, 5.3, and 5.5 prove that
every listed exception has minimum dimension $2t+1$.

For any other realizable profile, apply a Fano automorphism so that it lies in
$C_1$. Lemma 6.2 says it is sharp. Therefore it is realized in dimension $2t$.
Undoing the relabelling preserves dimension and proves (5.3) for every
realizable profile.
\end{proof}


\section{Computational proof boundary}


The two main theorems combine ordinary proofs with finite exact certificates.
For clarity, the logical division is as follows.

\subsection*{General arguments independent of computation}


\begin{itemize}
\item rank subadditivity and the triangle inequalities;
\item nonrealizability of primitive line jumps by four-dimensional descent and
\end{itemize}

  Pfaffian parity;
\begin{itemize}
\item the exact singleton-face semigroup by commuting self-adjoint idempotents;
\item the lower bounds for the height-one and height-two exceptions;
\item the arbitrary-dimensional rank-two Pfaffian update identity;
\item nonsharpness of every $E_t$ and its propagation by Fano cuts;
\item orthogonal-sum propagation once a finite witness is available.

\end{itemize}

\subsection*{Finite exact certificates}


\begin{itemize}
\item completeness of the fourteen extreme rays and the 78-element Hilbert basis;
\item realization of the finite atom and coupling libraries;
\item completeness of the 30 primitive rays of the anchored chamber;
\item the 81-cone triangulation and its 226 residue classes;
\item exact equality of the conductor's 812 required bases with the frozen list;
\item sharp witness coverage of all 812 bases;
\item explicit low-dimensional witnesses used by the conductor.

\end{itemize}

The core reproduction command is

\begin{quote}\small\raggedright
\texttt{pwsh -NoProfile -File research/\allowbreak{}fano-half-rank-profiles/\allowbreak{}verification/\allowbreak{}run\_\allowbreak{}all.ps1}\par
\end{quote}


It uses Python with SymPy and a C++17 compiler and writes binaries only below
the operating system's temporary directory. The source hashes are recorded in
\nolinkurl{verification/SHA256SUMS}.

The bounded check with $\max h\le6$ and the fixed-seed witness searches are
regressions and construction tools. They are not substituted for the general
nonrealizability arguments.

\section{Relation to prior work and limits of the result}


Pott, Schmidt, and Zhou \cite{pott2016} count pairs of alternating matrices according to the
three ranks of $(A,B,A+B)$. This is the two-dimensional, three-point precursor
of the present seven-point profile. Their result does not give the compatibility
semigroup or minimum effective dimension for a three-dimensional parameter
space.

Eisenbud and Koh \cite{eisenbud1994} classify three-dimensional spaces of $4\times4$ and
$5\times5$ alternating matrices over an algebraically closed field. Gow \cite{gow2017}
studies rank-related dimension bounds for subspaces of bilinear forms, and Qiao
\cite{qiao2021} enumerates alternating-matrix spaces over finite fields. The present result
does not refine those works as an orbit classification. Instead, it forgets the
isomorphism type and retains the seven labelled ranks, allowing arbitrary
ambient dimension over $\mathbb F_2$.

The following statements are outside the scope of the paper.

\begin{enumerate}
\item We do not classify alternating nets up to equivalence or count the nets with
   a given profile.
\item We do not treat arbitrary finite fields or odd characteristic.
\item We do not classify affine linear terms, Walsh signs, common zero-fibre
   histograms, or supports of three-dimensional subcodes of second-order
   Reed--Muller codes.
\item We do not infer formal peer-review status or global priority from the current
   literature search. The comparison above records the nearest public sources
   checked for this draft.

\end{enumerate}

\section{Conclusion}


The seven half-ranks of a three-parameter alternating net over $\mathbb F_2$
are governed by a nearly saturated Fano triangle semigroup. Its failure of
saturation has exactly two forms: primitive line jumps and singleton-face parity
cosets. After realizability is settled, the minimum effective dimension is even
and sharp except for a complete list of finite low-height orbits and one
cut-generated infinite family.

The result illustrates a useful division of labor. Pfaffian and idempotent
arguments explain the infinite obstructions, while finite exact polyhedral
certificates close the additive sufficiency problem. Neither bounded enumeration
nor unsuccessful search is used to assert a general theorem.

\begin{thebibliography}{9}
\footnotesize
\raggedright

\bibitem{pott2016}
A.~Pott, K.-U.~Schmidt, and Y.~Zhou,
\emph{Pairs of quadratic forms over finite fields},
Electron. J. Combin. 23(2) (2016), Paper 2.8.
\href{https://doi.org/10.37236/4072}{doi:10.37236/4072}.

\bibitem{eisenbud1994}
D.~Eisenbud and J.~Koh,
\emph{Nets of alternating matrices and the linear syzygy conjectures},
Adv. Math. 106 (1994), 1--35.
\href{https://doi.org/10.1006/aima.1994.1047}{doi:10.1006/aima.1994.1047}.

\bibitem{gow2017}
R.~Gow,
\emph{Rank-related dimension bounds for subspaces of bilinear forms over finite fields},
arXiv:1703.07266 (2017).
\href{https://arxiv.org/abs/1703.07266}{arXiv:1703.07266}.

\bibitem{qiao2021}
Y.~Qiao,
\emph{Enumerating alternating matrix spaces over finite fields with explicit coordinates},
Discrete Math. 344 (2021), 112580.
\href{https://doi.org/10.1016/j.disc.2021.112580}{doi:10.1016/j.disc.2021.112580}.

\end{thebibliography}

\section*{AI-assisted research and writing disclosure}

The author used OpenAI Codex for exploratory computation, proof checking, and
language editing. The author reviewed the mathematical claims and assumes
responsibility for the manuscript.

\end{document}
